\chapter{Compiling and Running Mathilda}
\label{ch:building}

Before you can compute anything, you have to get Mathilda onto your machine and
build it. For most newcomers this is the single largest hurdle---not the
mathematics, but the toolchain\index{building!prerequisites}---so this chapter treats it carefully and
completely. By the end you will have obtained the source, installed the handful of
libraries it needs, compiled the program, started an interactive session, and
learned your way around the read--eval--print loop. If you have built open-source C
programs before, much of this will be familiar and you can skim; if you have not,
follow it line by line and you will be fine.

Mathilda is distributed as \emph{source code}, not as a pre-built binary. That is a
deliberate consequence of the mission of Chapter~\ref{ch:about}: a system meant to
be read, checked, and carried forward is one you compile yourself, on your own
hardware, from source you can inspect. The good news is that the build is ordinary
and undramatic---a C project with a plain \texttt{makefile}, a few well-known
libraries, and no exotic build system to learn.

\section{Obtaining the source}
\label{sec:obtaining}
\index{building!obtaining the source}

Mathilda lives in a public Git repository. To get a copy, clone it and step
inside:

\begin{shell}
$ git clone https://github.com/stblake/Mathilda.git
$ cd Mathilda
\end{shell}

\noindent
Everything the rest of this chapter does happens from inside that top-level
\texttt{Mathilda/} directory. Two files there are worth knowing about from the
start: \texttt{README.md} is the quick-start and the authoritative list of
dependencies, and \texttt{SPEC.md} describes the system's architecture. You do not
need to read either to build the program, but they are where to turn when something
is unusual about your platform.

\section{The toolchain and its dependencies}
\label{sec:prerequisites}

Mathilda is written in C99 and leans on a small set of best-in-class numerical
libraries (introduced in \S\ref{sec:libraries}). Building it therefore needs a C
compiler and those libraries' development headers. Only three of the libraries are
required; the rest are optional and the build \emph{degrades gracefully}\index{graceful degradation} without
them---it detects what is present, switches on the corresponding features, and
switches off the ones whose libraries are missing, always producing a working
binary.

\subsection{A real GCC, not clang}
\label{sec:gcc-not-clang}
\index{GCC!required, not clang}

The one platform rule that trips people up: Mathilda must be built with a genuine
\textbf{GCC}, never Apple's or LLVM's \texttt{clang}. The build enforces this for
you. Its makefile auto-detects a real GCC---trying \texttt{gcc-16}, \texttt{gcc-15},
\ldots\ down to a plain \texttt{gcc}---and if the compiler it is about to use reports
itself as clang, it stops immediately with a message telling you how to fix it,
rather than silently producing an unsupported binary.%
\calloutnote{Under the hood}{The rule is not arbitrary. On macOS the plain
\texttt{gcc} command is a symlink to Apple's clang, which both miscompiles a few of
Mathilda's constructs and---more importantly---hides the strict-C99 portability
warnings that the project's Linux continuous-integration build treats as errors.
Building with real GCC on every platform keeps one standard of correctness. The fix
is always the same: install GCC from Homebrew.}
On Linux the plain \texttt{gcc} is already real GCC, so there is nothing to do. On
macOS, install one from Homebrew (see below); the makefile will then find and prefer
it automatically. If you ever need to point it at a specific compiler, pass it
explicitly, for example \texttt{make CC=gcc-16}.

\subsection{The libraries}
\label{sec:libraries-list}

\paragraph{Required.} These three must be present or the build cannot proceed:
\begin{itemize}
  \item \textbf{GMP}\index{GMP}, the GNU Multiple Precision Arithmetic Library---arbitrary-precision
        integers, the foundation of all of Mathilda's exact arithmetic.
  \item \textbf{MPFR}\index{MPFR}---correctly rounded arbitrary-precision reals, so that \B{N}
        can give you a hundred trustworthy digits. It is on by default; you can build
        without it, but almost no one should.
  \item \textbf{GNU~Readline}\index{GNU Readline}---the line editing, history, and command recall of the
        interactive session.
\end{itemize}

\paragraph{Optional.} Each of these unlocks a family of features when present and is
simply omitted when absent:
\begin{itemize}
  \item \textbf{FLINT}\index{FLINT} (version~3.0 or newer)---fast, rigorous polynomial arithmetic
        over algebraic number fields and rigorous numerics for many special
        functions. Distributions that ship only FLINT~2.x are detected as too old
        and the accelerated paths are switched off automatically.
  \item \textbf{GMP-ECM}\index{GMP-ECM}---the Elliptic Curve Method for the harder cases of integer
        factorization.
  \item \textbf{LAPACK}\index{LAPACK} and \textbf{BLAS}---fast machine-precision linear algebra. On
        macOS the build finds Apple's \textbf{Accelerate} framework and you need
        install nothing.
  \item \textbf{FFTW}---fast Fourier transforms; without it \B{Fourier} falls back to
        a slower $O(n^2)$ transform.
  \item \textbf{Raylib}\index{Raylib}---the interactive window that \B{Plot}, \B{Plot3D}, and \B{Show}
        draw into. Without it those functions print a text placeholder and return
        normally (see Chapter~\ref{ch:graphics}).
\end{itemize}

\noindent
POSIX threads, used to spread the packed-array kernels across your cores, are on by
default on macOS and Linux and need no extra package. \textbf{CMake} is needed only
to build the test suite (\S\ref{sec:tests}).

Every optional backend is governed by a build-time flag you can force on or off. The
defaults auto-detect, so you rarely need to touch them, but they are worth knowing:

\begin{table}[htbp]
  \centering
  \small
  \begin{tabular}{@{}llp{0.52\linewidth}@{}}
    \toprule
    \textbf{Flag} & \textbf{Default} & \textbf{What it enables} \\
    \midrule
    \mcode{USE\_MPFR}     & \mcode{1} & Arbitrary-precision reals: \mcode{N[expr, prec]}, precision tracking. \\
    \mcode{USE\_FLINT}    & \mcode{1} & Fast FLINT ($\geq 3.0$) kernels for number fields and rigorous numerics. \\
    \mcode{USE\_LAPACK}   & \mcode{1} & Fast machine-precision linear algebra (Accelerate on macOS). \\
    \mcode{USE\_ECM}      & \mcode{1} & Elliptic Curve Method factorization via GMP-ECM. \\
    \mcode{USE\_FFTW}     & \mcode{1} & FFTW-backed \mcode{Fourier}; naive $O(n^2)$ transform when off. \\
    \mcode{USE\_THREADS}  & \mcode{1} & POSIX-thread parallelism for the packed-array kernels. \\
    \mcode{USE\_GRAPHICS} & \mcode{1} & Interactive 2D/3D plot windows via Raylib. \\
    \bottomrule
  \end{tabular}
  \caption{Optional-backend flags. Each auto-detects; append e.g.\ \mcode{USE\_LAPACK=0}
  to a \mcode{make} command to force one off.}
  \label{tab:use-flags}
\end{table}

\subsection{Installing the dependencies}
\label{sec:installing-deps}
\index{building!installing dependencies}

On \textbf{Debian or Ubuntu}, install the three required libraries and as many of the
optional ones as you want:

\begin{shell}
# Required
$ sudo apt install libgmp-dev libmpfr-dev libreadline-dev

# Optional, each auto-detected by the build
$ sudo apt install libflint-dev          # FLINT >= 3.0 (Ubuntu 24.04+/Debian 12+)
$ sudo apt install libecm-dev            # GMP-ECM factorization
$ sudo apt install liblapacke-dev libopenblas-dev   # LAPACK/BLAS
$ sudo apt install libfftw3-dev          # FFTW
$ sudo apt install libraylib-dev         # interactive graphics
$ sudo apt install cmake                 # only for the test suite
\end{shell}

\noindent
On \textbf{Fedora or RHEL} the package names are \texttt{gmp-devel},
\texttt{mpfr-devel}, \texttt{readline-devel}, and then \texttt{flint-devel},
\texttt{gmp-ecm-devel}, \texttt{lapack-devel}/\texttt{openblas-devel},
\texttt{fftw-devel}, and \texttt{cmake}.

On \textbf{macOS} with Homebrew, install a real GCC (\S\ref{sec:gcc-not-clang}) and
the libraries; Accelerate is already part of the system, so you do not install
LAPACK:

\begin{shell}
$ brew install gcc gmp mpfr readline cmake
# Optional
$ brew install flint fftw raylib gmp-ecm
\end{shell}

\section{Building}
\label{sec:building}
\index{building!compiling}

With the dependencies in place, build the program with a single command. The
makefile auto-discovers every source file, configures the internal dependencies, and
links the executable with \texttt{-std=c99 -O3}:

\begin{shell}
$ make -j
\end{shell}

\noindent
The \texttt{-j} flag builds in parallel across your cores; on Linux you can be
explicit with \texttt{make -j\$(nproc)}. The first build compiles several hundred
source files and takes a little while; afterwards, \texttt{make} rebuilds only what
changed. When it finishes you will have an executable called \texttt{Mathilda} in the
current directory.

Want a leaner binary, or missing a library you do not care about? Turn any optional
backend off explicitly (Table~\ref{tab:use-flags}). For instance, to build without
MPFR and without the LAPACK acceleration:

\begin{shell}
$ make -j USE_MPFR=0 USE_LAPACK=0
\end{shell}

\paragraph{If the build stops.} The three failures new users hit, and their fixes:
\begin{itemize}
  \item \emph{``\,'gcc' is clang, not GCC\,''}---you are on macOS with only Apple's
        compiler. Run \texttt{brew install gcc} and build again; the makefile will
        find the real GCC on its own (\S\ref{sec:gcc-not-clang}).
  \item \emph{A missing header such as \texttt{gmp.h} or \texttt{readline/readline.h}}---%
        a required library's development package is not installed. Install it from
        \S\ref{sec:installing-deps} and rebuild.
  \item \emph{FLINT features quietly absent}---your distribution ships FLINT~2.x.
        The build detects this, falls back to \texttt{USE\_FLINT=0}, and still
        succeeds; install FLINT~$\geq 3.0$ from source or a backport if you want the
        accelerated paths.
\end{itemize}

\section{Your first session}
\label{sec:first-session}
\index{REPL!starting}

Start the interactive session by running the binary with no arguments:

\begin{shell}
$ ./Mathilda

Mathilda <version> - A small, open source computer algebra system.

This program is free, open source software and comes with ABSOLUTELY NO WARRANTY.

Press Return to evaluate. An open bracket, string or comment continues
on the next line; press Esc then Return to force evaluation.
Exit by evaluating Quit[] or CONTROL-C.

In[1]:=
\end{shell}

\noindent
The banner's first line names the version you just built and the libraries it was
compiled against; the same information is available at any time from
\B{\$Version}\index{version!\$Version}, and the version number alone from \B{\$VersionNumber}:

\begin{codepairs}
\pair{02-building/version/1}{The full banner: Mathilda's version and every library
it was built against---exactly the build you are running.}
\pair{02-building/version/2}{Just the number, as a machine-readable value you can
compare against in a script.}
\end{codepairs}

The \texttt{In[1]:=} is the \emph{prompt}\index{REPL!In[] and Out[]}: Mathilda is waiting for you to type an
expression. Type one, press Return, and it prints the result on an
\texttt{Out[1]=} line and prompts again. Here is the whole of a first session---exact
arithmetic, an exact symbolic value, and a machine-precision number on demand:

\begin{codepairs}
\pair{02-building/first-session/1}{Ordinary arithmetic. The answer is the exact
integer~$4$.}
\pair{02-building/first-session/2}{\B{Sin} of $\pi/6$ is returned \emph{exactly}, as
the rational $1/2$---not a decimal.}
\pair{02-building/first-session/3}{Ask for a number with \B{N} and you get one: here
a quick machine-precision value of $\pi$.}
\end{codepairs}

\noindent
That is the entire rhythm of the system: you type an expression, Mathilda evaluates
it to a fixed point, and it shows you the result. Everything else in this book is
variations on that loop.

\section{How the REPL works}
\label{sec:repl}
\index{REPL}

The interactive loop---read an expression, evaluate it, print it, repeat---is where
you will spend most of your time, so it is worth knowing its conveniences.

\subsection{Results persist, and you can name them}
\label{sec:session-state}

A session has \emph{memory}. Everything you define stays in force until you leave, so
you can name a result on one line and use it on the next:

\begin{codepairs}
\pair{02-building/session-state/1}{Assign with \mcode{=}. The assignment returns the
value, so you see it echoed.}
\pair{02-building/session-state/2}{\mcode{x} now stands for $1024$ everywhere.}
\pair{02-building/session-state/3}{And keeps standing for it until you reassign or
clear it.}
\end{codepairs}

Beyond your own names, the session numbers everything for you.
Each result is stored as \mcode{Out[n]}, matching the \texttt{Out[n]=} label beside
it, and each input as \mcode{In[n]}. The special symbol \mcode{\%} is shorthand for
the most recent result and \mcode{\%\%} for the one before it, so in an interactive
session \mcode{\% + 1} adds one to whatever was just printed, and \mcode{Out[3]}
recalls the third result exactly.%
\calloutnote{Under the hood}{\mcode{In[n]}, \mcode{Out[n]}, and \mcode{\%} are a
feature of the \emph{interactive} loop---the REPL records each result as it prints
it. A script run with \mcode{-file} (\S\ref{sec:scripts}) and the machine pipe used
by front-ends do not keep this history, so reach for a named variable, as above, for
anything you want to reuse in a script.}

A trailing semicolon suppresses the printed output. The expression is still
evaluated and its value still stored---you just do not see it, which is exactly what
you want for setup lines and large intermediate results:

\begin{codepairs}
\pair{02-building/suppression/1}{The semicolon suppresses the \texttt{Out[]} line;
the assignment still happens.}
\pair{02-building/suppression/2}{Proof that it did: \mcode{a} is $3$, so \mcode{a\^{}2}
is $9$.}
\end{codepairs}

\subsection{Editing, history, and multiline input}
\label{sec:editing}
\index{REPL!line editing and history}

Because the session is built on GNU~Readline, it behaves like a modern shell prompt.
You can move along the line with the arrow keys, edit in place, and---most
usefully---press the up arrow to recall previous inputs and re-run or tweak them. The
Emacs-style key bindings work too: \texttt{Ctrl-A} to the start of the line,
\texttt{Ctrl-E} to the end, \texttt{Ctrl-R} to search backwards through history.

For an expression too long or too structured to sit comfortably on one line, you need
do nothing special: just keep typing.\index{REPL!multiline input} Pressing Return
evaluates the buffer \emph{only} once it forms a complete
expression---every \mcode{(}, \mcode{[} and \mcode{\{} closed, and every string and
comment finished. Until then, Return simply opens a fresh line and waits, so a
bracket you have opened but not yet closed carries you onto the next line
automatically:

\begin{shell}
In[1]:= Table[i^2,
{i, 1, 4}]
Out[1]= {1, 4, 9, 16}
\end{shell}

\noindent
You close the final brace, press Return, and the whole thing runs as one
expression---the line breaks inside it are just whitespace, exactly as they are in a
script. Should you ever want to submit an \emph{incomplete} line as it stands---to
see the parser's error, say, rather than being carried onto a new line
forever---press \texttt{Esc} then \texttt{Return} (or \texttt{Alt-Return}), which
forces evaluation regardless.%
\calloutnote{Under the hood}{This in-place multiline editing needs a genuine
GNU~Readline. macOS ships a \texttt{libedit} stand-in under the same header name that
cannot do it, so the build prefers a Homebrew \texttt{readline} when one is present
and otherwise falls back to plain one-line-at-a-time entry.}

\subsection{Getting help}
\label{sec:help}
\index{REPL!getting help}

To see what a function does, type a question mark before its name---\mcode{?Sin},
say---and Mathilda prints its usage message: a one-line signature and a short
description, drawn from the function's own built-in documentation. Every builtin
carries one. This is the same text that titles the function's reference card; for
\B{Sin} it reads:

\usagebox{Sin}

To \emph{discover} functions rather than look one up, \B{Names} lists every symbol
matching a string pattern---a quick way to see what a fresh session already knows:

\begin{codepairs}
\pair{02-building/help/1}{Every built-in name beginning with \mcode{Bessel}. The
\mcode{*} is a wildcard.}
\end{codepairs}

\noindent
For the fuller machine-readable story about a symbol---its attributes and
definitions as well as its usage---there is also \B{Information}, of which
\mcode{?name} is the everyday shorthand.

\subsection{Leaving the session}
\label{sec:leaving}

To quit, evaluate \mcode{Quit[]} or press \texttt{Ctrl-D} (end-of-input) at the
prompt. \texttt{Ctrl-C} interrupts as well. Any of them returns you to your shell.

\section{Running scripts and other modes}
\label{sec:scripts}
\index{building!running scripts}

The REPL is one of three ways Mathilda can run. It is chosen automatically: if you
name a file, Mathilda runs it as a script; otherwise, if its input is an interactive
terminal, it starts the REPL you have just met; and if its input is a pipe, it speaks
a machine protocol instead (below).

\paragraph{Scripts.} Point Mathilda at a file of expressions with \texttt{-file} (a
bare filename means the same thing) and it evaluates each one in order and exits.
Crucially, \emph{nothing is echoed}: there are no \texttt{In[]}/\texttt{Out[]} lines,
so the output is exactly and only what the script chooses to \B{Print}. Given a file
\texttt{area.m}:

\begin{shell}
$ cat area.m
(* area.m -- a tiny script. Nothing is echoed; only what we Print appears. *)
r = 3;
Print["Area of a disk of radius ", r, " is ", N[Pi r^2]]

$ ./Mathilda -file area.m
Area of a disk of radius 3 is 28.2743
\end{shell}

\noindent
This is the mode to use for batch work, for reproducible computations kept under
version control, and for anything you want to pipe into another program. (It is also
how this very book is built: every transcript you see was produced by feeding its
inputs to the real binary.)

\paragraph{The command line.} A few options round out the picture; \texttt{--help}
lists them all:

\begin{shell}
$ ./Mathilda --help
Mathilda <version> - a small, open source computer algebra system.

Usage: ./Mathilda [options] [file]

  -file <path>    evaluate every expression in <path>, then exit
  -h, --help      show this message and exit
  -v, --version   print version information and exit

A bare <file> argument is equivalent to -file <file>. With no file,
Mathilda starts the interactive REPL when stdin is a terminal, and
otherwise speaks the NDJSON pipe protocol on stdio.
\end{shell}

\paragraph{The pipe protocol.} That last sentence describes the third mode. When
Mathilda's input is not a terminal and no script was named, it speaks a line-oriented
JSON protocol on its standard input and output, designed for a front-end---a
notebook, an editor plugin, a web interface---to drive it programmatically.%
\calloutnote{Under the hood}{The protocol is \emph{NDJSON}: one JSON request per
line in, one or more JSON responses per line out, each result carrying both a text
form and a \LaTeX{} form so a graphical front-end can render it. You will not use
this mode by hand, but it is what lets Mathilda sit behind a richer interface than a
terminal.}

\section{Building and running the tests}
\label{sec:tests}
\index{building!test suite}

Mathilda ships a large suite of C-based unit tests. You do not need them to use the
system, but running them is the surest way to confirm your build is sound---and if
you go on to modify Mathilda (Chapter~\ref{ch:contributing}), they are how you check
you have not broken anything. The suite is built with CMake, which auto-detects the
same optional backends as the main build:

\begin{shell}
$ cd tests
$ mkdir -p build && cd build
$ cmake ..
$ make -j

# Run every test binary
$ for t in *_tests; do ./$t; done
\end{shell}

\bigskip
\noindent
That is everything you need to obtain, build, and run Mathilda, and to find your way
around a session. With a working binary in front of you and the prompt waiting, we
can turn to the mathematics. The next chapter is a whirlwind tour of what the system
can do---arithmetic, number theory, algebra, calculus, and more---so that the
detailed chapters afterwards feel like a return to places you have already seen.
